% Toolbox Vignette of the Spatial Econometrics Toolbox
% Compile with: julia -e 'using PollisTeX; display(compile("spatial-econometrics.tex"))'
% The look of a Pollis Toolbox Vignette (the fonts and colours of the Pollis New Methods documents).
% Set \vignetteauthor, \vignettedate and, if not CC BY-SA 4.0, \vignettelicense after % The look of a Pollis Toolbox Vignette (the fonts and colours of the Pollis New Methods documents).
% Set \vignetteauthor, \vignettedate and, if not CC BY-SA 4.0, \vignettelicense after % The look of a Pollis Toolbox Vignette (the fonts and colours of the Pollis New Methods documents).
% Set \vignetteauthor, \vignettedate and, if not CC BY-SA 4.0, \vignettelicense after \input{vignette_style}.
\documentclass[10pt,a4paper]{article}

\usepackage[a4paper,left=18mm,right=18mm,top=15mm,bottom=17mm,footskip=8mm]{geometry}
\usepackage{iftex}
\ifPDFTeX\usepackage[T1]{fontenc}\fi
\usepackage[default,semibold]{sourcesanspro}
\usepackage{textcomp}
\usepackage[italic]{mathastext}
\usepackage{amsmath}
\usepackage{microtype}
\usepackage[table]{xcolor}
\usepackage{tikz}
\usepackage[most]{tcolorbox}
\usepackage{fontawesome5}
\usepackage{tabularx,booktabs}
\usepackage{enumitem}
\usepackage{fancyhdr}
\usepackage[hidelinks]{hyperref}

% ---------------------------------------------------------------- colours (as in the brochure)
\definecolor{navy}{HTML}{1C3A5E}
\definecolor{pblue}{HTML}{3B78B5}
\definecolor{amber}{HTML}{F2B33D}
\definecolor{deepamber}{HTML}{D9902A}
\definecolor{pale}{HTML}{EEF3F9}
\definecolor{ink}{HTML}{2A2F36}
\definecolor{logonavy}{HTML}{173D63}
\definecolor{logodark}{HTML}{C98A1E}
\definecolor{logoamber}{HTML}{F2B446}
\definecolor{logocream}{HTML}{FEEFD2}
\definecolor{soft}{HTML}{5E6773}
\definecolor{agree}{HTML}{2F7D4F}
\definecolor{differ}{HTML}{B23B3B}
\color{ink}

\setlength{\parindent}{0pt}
\setlength{\parskip}{4pt}
\setlist{nosep,leftmargin=*,itemsep=2pt}
\renewcommand{\arraystretch}{1.15}

\pagestyle{fancy}
\fancyhf{}
\renewcommand{\headrulewidth}{0pt}
\fancyfoot[L]{\footnotesize\color{soft}\textbf{Pollis}\enspace Toolbox Vignette\enspace·\enspace\nmtheme}
\fancyfoot[R]{\footnotesize\color{soft}\thepage}
\newcommand{\nmtheme}{}
\newcommand{\vignetteauthor}{}
\newcommand{\vignettedate}{\today}
\newcommand{\vignettelicense}{\href{https://creativecommons.org/licenses/by-sa/4.0/}{CC BY-SA 4.0}}

% The Pollis logo, redrawn so it prints sharply at any size
\newcommand{\pollislogo}[1]{%
  \begin{tikzpicture}[x=#1/512,y=-#1/512,baseline=(current bounding box.center)]
    \fill[logonavy] (0,0) rectangle (512,512);
    \fill[logodark] (36,422) -- (104,256) -- (172,422) -- cycle;
    \fill[logoamber] (104,422) -- (203,180) -- (302,422) -- cycle;
    \fill[logocream] (203,422) -- (339,89) -- (476,422) -- cycle;
  \end{tikzpicture}}

% Top of the first page: logo, name; then the kicker, title and a one-line promise
\newcommand{\nmhead}[3]{%
  \renewcommand{\nmtheme}{#1}%
  \thispagestyle{fancy}%
  \begin{minipage}[c]{13mm}\pollislogo{11mm}\end{minipage}%
  \begin{minipage}[c]{\dimexpr\linewidth-13mm\relax}
    {\color{navy}\Large\bfseries Pollis}\enspace{\color{soft}Toolbox Vignette}\hfill
    {\color{deepamber}\small\bfseries\MakeUppercase{#1}}
  \end{minipage}\par\vspace{7mm}
  {\color{navy}\fontsize{22}{26}\selectfont\bfseries #2}\par\vspace{4pt}
  {\color{soft}\large #3}\par\vspace{5mm}}

\newcommand{\nmsection}[1]{\par\vspace{5pt}{\color{navy}\large\bfseries #1}\par\vspace{1pt}}
\newcommand{\nmsub}[1]{\par\vspace{2pt}{\color{deepamber}\bfseries #1}\par}
\newcommand{\gt}{\,{\color{deepamber}\bfseries ›}\,}
\newcommand{\menupath}[1]{{\bfseries\color{navy}#1}}
\newcommand{\fn}[1]{\texttt{\small #1}}
% Citations: \nmcite{key}{Breiman (2001)} links to the entry \nmref{key}{...} in the References
\newcommand{\nmcite}[2]{\hyperlink{ref:#1}{\color{pblue}#2}}
\newcommand{\nmref}[2]{\raisebox{1.1em}[0pt][0pt]{\hypertarget{ref:#1}{}}#2\par}
\newcommand{\nmdoi}[1]{\href{https://doi.org/#1}{\color{pblue}doi:#1}}
% Where the package is: its name and a link (its GitHub repository, once published)
\newcommand{\package}[2]{\href{#2}{\fn{#1.jl}}}

% Key facts at the top: source, missing in Julia, package, result
\newcommand{\factsheet}[4]{%
  \begin{tcolorbox}[enhanced, colback=pale, colframe=pale, boxrule=0pt, sharp corners,
    left=7pt, right=7pt, top=5pt, bottom=5pt, fontupper=\small]
  \begin{tabularx}{\linewidth}{@{}>{\leavevmode\bfseries\color{navy}}p{27mm}>{\leavevmode}X@{}}
  Tested against & #1\\
  Missing in Julia & #2\\
  Package & #3\\
  Result & #4
  \end{tabularx}
  \end{tcolorbox}}

% Boxes
\newtcolorbox{pbox}[2][]{enhanced, colback=pale, colframe=amber, boxrule=0pt,
  leftrule=3pt, sharp corners, left=7pt, right=7pt, top=6pt, bottom=6pt,
  fontupper=\small\raggedright, before upper={{\color{navy}\normalsize\bfseries #2}\par\smallskip}, #1}
\newtcolorbox{nbox}[2][]{enhanced, colback=navy, colframe=navy, boxrule=0pt,
  sharp corners, left=8pt, right=8pt, top=7pt, bottom=7pt, coltext=white,
  fontupper=\small\raggedright, before upper={{\color{amber}\normalsize\bfseries #2}\par\smallskip}, #1}

% Value tables: published beside Julia
\newenvironment{valuetable}[1]{%
  \par\small\tabularx{\linewidth}{#1}\toprule}{\bottomrule\endtabularx\par}
\newcommand{\ok}{{\color{agree}\faCheck}}
\newcommand{\no}{{\color{differ}\faTimes}}

% End of the document: credit and licence
\newcommand{\nmcredit}{%
  \par\vfill
  \begin{minipage}[b]{0.78\linewidth}\footnotesize\color{soft}
  A Pollis Toolbox Vignette, written from the panels of the toolbox.\par
  \textcopyright\ \vignetteauthor\enspace·\enspace Licensed under \vignettelicense\enspace·\enspace\vignettedate
  \end{minipage}\hfill
  \begin{minipage}[b]{0.18\linewidth}\raggedleft\pollislogo{9mm}\end{minipage}}
.
\documentclass[10pt,a4paper]{article}

\usepackage[a4paper,left=18mm,right=18mm,top=15mm,bottom=17mm,footskip=8mm]{geometry}
\usepackage{iftex}
\ifPDFTeX\usepackage[T1]{fontenc}\fi
\usepackage[default,semibold]{sourcesanspro}
\usepackage{textcomp}
\usepackage[italic]{mathastext}
\usepackage{amsmath}
\usepackage{microtype}
\usepackage[table]{xcolor}
\usepackage{tikz}
\usepackage[most]{tcolorbox}
\usepackage{fontawesome5}
\usepackage{tabularx,booktabs}
\usepackage{enumitem}
\usepackage{fancyhdr}
\usepackage[hidelinks]{hyperref}

% ---------------------------------------------------------------- colours (as in the brochure)
\definecolor{navy}{HTML}{1C3A5E}
\definecolor{pblue}{HTML}{3B78B5}
\definecolor{amber}{HTML}{F2B33D}
\definecolor{deepamber}{HTML}{D9902A}
\definecolor{pale}{HTML}{EEF3F9}
\definecolor{ink}{HTML}{2A2F36}
\definecolor{logonavy}{HTML}{173D63}
\definecolor{logodark}{HTML}{C98A1E}
\definecolor{logoamber}{HTML}{F2B446}
\definecolor{logocream}{HTML}{FEEFD2}
\definecolor{soft}{HTML}{5E6773}
\definecolor{agree}{HTML}{2F7D4F}
\definecolor{differ}{HTML}{B23B3B}
\color{ink}

\setlength{\parindent}{0pt}
\setlength{\parskip}{4pt}
\setlist{nosep,leftmargin=*,itemsep=2pt}
\renewcommand{\arraystretch}{1.15}

\pagestyle{fancy}
\fancyhf{}
\renewcommand{\headrulewidth}{0pt}
\fancyfoot[L]{\footnotesize\color{soft}\textbf{Pollis}\enspace Toolbox Vignette\enspace·\enspace\nmtheme}
\fancyfoot[R]{\footnotesize\color{soft}\thepage}
\newcommand{\nmtheme}{}
\newcommand{\vignetteauthor}{}
\newcommand{\vignettedate}{\today}
\newcommand{\vignettelicense}{\href{https://creativecommons.org/licenses/by-sa/4.0/}{CC BY-SA 4.0}}

% The Pollis logo, redrawn so it prints sharply at any size
\newcommand{\pollislogo}[1]{%
  \begin{tikzpicture}[x=#1/512,y=-#1/512,baseline=(current bounding box.center)]
    \fill[logonavy] (0,0) rectangle (512,512);
    \fill[logodark] (36,422) -- (104,256) -- (172,422) -- cycle;
    \fill[logoamber] (104,422) -- (203,180) -- (302,422) -- cycle;
    \fill[logocream] (203,422) -- (339,89) -- (476,422) -- cycle;
  \end{tikzpicture}}

% Top of the first page: logo, name; then the kicker, title and a one-line promise
\newcommand{\nmhead}[3]{%
  \renewcommand{\nmtheme}{#1}%
  \thispagestyle{fancy}%
  \begin{minipage}[c]{13mm}\pollislogo{11mm}\end{minipage}%
  \begin{minipage}[c]{\dimexpr\linewidth-13mm\relax}
    {\color{navy}\Large\bfseries Pollis}\enspace{\color{soft}Toolbox Vignette}\hfill
    {\color{deepamber}\small\bfseries\MakeUppercase{#1}}
  \end{minipage}\par\vspace{7mm}
  {\color{navy}\fontsize{22}{26}\selectfont\bfseries #2}\par\vspace{4pt}
  {\color{soft}\large #3}\par\vspace{5mm}}

\newcommand{\nmsection}[1]{\par\vspace{5pt}{\color{navy}\large\bfseries #1}\par\vspace{1pt}}
\newcommand{\nmsub}[1]{\par\vspace{2pt}{\color{deepamber}\bfseries #1}\par}
\newcommand{\gt}{\,{\color{deepamber}\bfseries ›}\,}
\newcommand{\menupath}[1]{{\bfseries\color{navy}#1}}
\newcommand{\fn}[1]{\texttt{\small #1}}
% Citations: \nmcite{key}{Breiman (2001)} links to the entry \nmref{key}{...} in the References
\newcommand{\nmcite}[2]{\hyperlink{ref:#1}{\color{pblue}#2}}
\newcommand{\nmref}[2]{\raisebox{1.1em}[0pt][0pt]{\hypertarget{ref:#1}{}}#2\par}
\newcommand{\nmdoi}[1]{\href{https://doi.org/#1}{\color{pblue}doi:#1}}
% Where the package is: its name and a link (its GitHub repository, once published)
\newcommand{\package}[2]{\href{#2}{\fn{#1.jl}}}

% Key facts at the top: source, missing in Julia, package, result
\newcommand{\factsheet}[4]{%
  \begin{tcolorbox}[enhanced, colback=pale, colframe=pale, boxrule=0pt, sharp corners,
    left=7pt, right=7pt, top=5pt, bottom=5pt, fontupper=\small]
  \begin{tabularx}{\linewidth}{@{}>{\leavevmode\bfseries\color{navy}}p{27mm}>{\leavevmode}X@{}}
  Tested against & #1\\
  Missing in Julia & #2\\
  Package & #3\\
  Result & #4
  \end{tabularx}
  \end{tcolorbox}}

% Boxes
\newtcolorbox{pbox}[2][]{enhanced, colback=pale, colframe=amber, boxrule=0pt,
  leftrule=3pt, sharp corners, left=7pt, right=7pt, top=6pt, bottom=6pt,
  fontupper=\small\raggedright, before upper={{\color{navy}\normalsize\bfseries #2}\par\smallskip}, #1}
\newtcolorbox{nbox}[2][]{enhanced, colback=navy, colframe=navy, boxrule=0pt,
  sharp corners, left=8pt, right=8pt, top=7pt, bottom=7pt, coltext=white,
  fontupper=\small\raggedright, before upper={{\color{amber}\normalsize\bfseries #2}\par\smallskip}, #1}

% Value tables: published beside Julia
\newenvironment{valuetable}[1]{%
  \par\small\tabularx{\linewidth}{#1}\toprule}{\bottomrule\endtabularx\par}
\newcommand{\ok}{{\color{agree}\faCheck}}
\newcommand{\no}{{\color{differ}\faTimes}}

% End of the document: credit and licence
\newcommand{\nmcredit}{%
  \par\vfill
  \begin{minipage}[b]{0.78\linewidth}\footnotesize\color{soft}
  A Pollis Toolbox Vignette, written from the panels of the toolbox.\par
  \textcopyright\ \vignetteauthor\enspace·\enspace Licensed under \vignettelicense\enspace·\enspace\vignettedate
  \end{minipage}\hfill
  \begin{minipage}[b]{0.18\linewidth}\raggedleft\pollislogo{9mm}\end{minipage}}
.
\documentclass[10pt,a4paper]{article}

\usepackage[a4paper,left=18mm,right=18mm,top=15mm,bottom=17mm,footskip=8mm]{geometry}
\usepackage{iftex}
\ifPDFTeX\usepackage[T1]{fontenc}\fi
\usepackage[default,semibold]{sourcesanspro}
\usepackage{textcomp}
\usepackage[italic]{mathastext}
\usepackage{amsmath}
\usepackage{microtype}
\usepackage[table]{xcolor}
\usepackage{tikz}
\usepackage[most]{tcolorbox}
\usepackage{fontawesome5}
\usepackage{tabularx,booktabs}
\usepackage{enumitem}
\usepackage{fancyhdr}
\usepackage[hidelinks]{hyperref}

% ---------------------------------------------------------------- colours (as in the brochure)
\definecolor{navy}{HTML}{1C3A5E}
\definecolor{pblue}{HTML}{3B78B5}
\definecolor{amber}{HTML}{F2B33D}
\definecolor{deepamber}{HTML}{D9902A}
\definecolor{pale}{HTML}{EEF3F9}
\definecolor{ink}{HTML}{2A2F36}
\definecolor{logonavy}{HTML}{173D63}
\definecolor{logodark}{HTML}{C98A1E}
\definecolor{logoamber}{HTML}{F2B446}
\definecolor{logocream}{HTML}{FEEFD2}
\definecolor{soft}{HTML}{5E6773}
\definecolor{agree}{HTML}{2F7D4F}
\definecolor{differ}{HTML}{B23B3B}
\color{ink}

\setlength{\parindent}{0pt}
\setlength{\parskip}{4pt}
\setlist{nosep,leftmargin=*,itemsep=2pt}
\renewcommand{\arraystretch}{1.15}

\pagestyle{fancy}
\fancyhf{}
\renewcommand{\headrulewidth}{0pt}
\fancyfoot[L]{\footnotesize\color{soft}\textbf{Pollis}\enspace Toolbox Vignette\enspace·\enspace\nmtheme}
\fancyfoot[R]{\footnotesize\color{soft}\thepage}
\newcommand{\nmtheme}{}
\newcommand{\vignetteauthor}{}
\newcommand{\vignettedate}{\today}
\newcommand{\vignettelicense}{\href{https://creativecommons.org/licenses/by-sa/4.0/}{CC BY-SA 4.0}}

% The Pollis logo, redrawn so it prints sharply at any size
\newcommand{\pollislogo}[1]{%
  \begin{tikzpicture}[x=#1/512,y=-#1/512,baseline=(current bounding box.center)]
    \fill[logonavy] (0,0) rectangle (512,512);
    \fill[logodark] (36,422) -- (104,256) -- (172,422) -- cycle;
    \fill[logoamber] (104,422) -- (203,180) -- (302,422) -- cycle;
    \fill[logocream] (203,422) -- (339,89) -- (476,422) -- cycle;
  \end{tikzpicture}}

% Top of the first page: logo, name; then the kicker, title and a one-line promise
\newcommand{\nmhead}[3]{%
  \renewcommand{\nmtheme}{#1}%
  \thispagestyle{fancy}%
  \begin{minipage}[c]{13mm}\pollislogo{11mm}\end{minipage}%
  \begin{minipage}[c]{\dimexpr\linewidth-13mm\relax}
    {\color{navy}\Large\bfseries Pollis}\enspace{\color{soft}Toolbox Vignette}\hfill
    {\color{deepamber}\small\bfseries\MakeUppercase{#1}}
  \end{minipage}\par\vspace{7mm}
  {\color{navy}\fontsize{22}{26}\selectfont\bfseries #2}\par\vspace{4pt}
  {\color{soft}\large #3}\par\vspace{5mm}}

\newcommand{\nmsection}[1]{\par\vspace{5pt}{\color{navy}\large\bfseries #1}\par\vspace{1pt}}
\newcommand{\nmsub}[1]{\par\vspace{2pt}{\color{deepamber}\bfseries #1}\par}
\newcommand{\gt}{\,{\color{deepamber}\bfseries ›}\,}
\newcommand{\menupath}[1]{{\bfseries\color{navy}#1}}
\newcommand{\fn}[1]{\texttt{\small #1}}
% Citations: \nmcite{key}{Breiman (2001)} links to the entry \nmref{key}{...} in the References
\newcommand{\nmcite}[2]{\hyperlink{ref:#1}{\color{pblue}#2}}
\newcommand{\nmref}[2]{\raisebox{1.1em}[0pt][0pt]{\hypertarget{ref:#1}{}}#2\par}
\newcommand{\nmdoi}[1]{\href{https://doi.org/#1}{\color{pblue}doi:#1}}
% Where the package is: its name and a link (its GitHub repository, once published)
\newcommand{\package}[2]{\href{#2}{\fn{#1.jl}}}

% Key facts at the top: source, missing in Julia, package, result
\newcommand{\factsheet}[4]{%
  \begin{tcolorbox}[enhanced, colback=pale, colframe=pale, boxrule=0pt, sharp corners,
    left=7pt, right=7pt, top=5pt, bottom=5pt, fontupper=\small]
  \begin{tabularx}{\linewidth}{@{}>{\leavevmode\bfseries\color{navy}}p{27mm}>{\leavevmode}X@{}}
  Tested against & #1\\
  Missing in Julia & #2\\
  Package & #3\\
  Result & #4
  \end{tabularx}
  \end{tcolorbox}}

% Boxes
\newtcolorbox{pbox}[2][]{enhanced, colback=pale, colframe=amber, boxrule=0pt,
  leftrule=3pt, sharp corners, left=7pt, right=7pt, top=6pt, bottom=6pt,
  fontupper=\small\raggedright, before upper={{\color{navy}\normalsize\bfseries #2}\par\smallskip}, #1}
\newtcolorbox{nbox}[2][]{enhanced, colback=navy, colframe=navy, boxrule=0pt,
  sharp corners, left=8pt, right=8pt, top=7pt, bottom=7pt, coltext=white,
  fontupper=\small\raggedright, before upper={{\color{amber}\normalsize\bfseries #2}\par\smallskip}, #1}

% Value tables: published beside Julia
\newenvironment{valuetable}[1]{%
  \par\small\tabularx{\linewidth}{#1}\toprule}{\bottomrule\endtabularx\par}
\newcommand{\ok}{{\color{agree}\faCheck}}
\newcommand{\no}{{\color{differ}\faTimes}}

% End of the document: credit and licence
\newcommand{\nmcredit}{%
  \par\vfill
  \begin{minipage}[b]{0.78\linewidth}\footnotesize\color{soft}
  A Pollis Toolbox Vignette, written from the panels of the toolbox.\par
  \textcopyright\ \vignetteauthor\enspace·\enspace Licensed under \vignettelicense\enspace·\enspace\vignettedate
  \end{minipage}\hfill
  \begin{minipage}[b]{0.18\linewidth}\raggedleft\pollislogo{9mm}\end{minipage}}

\renewcommand{\vignetteauthor}{Pollis}
\renewcommand{\vignettedate}{19 November 2026}
\renewcommand{\vignettelicense}{AGPL-3.0}

\begin{document}

\nmhead{Spatial Econometrics Toolbox}
  {Spatial Econometrics Toolbox}
  {What its panels are for, and when to use each.}

Installing the toolbox adds the submenu \menupath{Extend \gt Spatial Econometrics}, with these panels. Each has an example you can edit and run, and actions on its results. All of them run on \package{SpatialEconometrics}{https://github.com/randsley/SpatialEconometrics.jl} and read area files, such as GeoPackage, GeoJSON or Shapefile, through the spatial extension of \package{DuckDB}{https://github.com/duckdb/duckdb}.

\medskip
{\small
\begin{tabularx}{\linewidth}{@{}>{\raggedright\arraybackslash}p{56mm}X@{}}
\toprule
\textbf{Panel} & \textbf{Use it when} \\\midrule
Spatial Lag and Error Models & An area's outcome depends on its neighbours' outcomes, their characteristics or shocks they share, in one cross-section. \\
Geographically Weighted Regression & A relationship seems to change from one part of the map to another. \\
Spatial Discrete Choice & The outcome is a yes or no, or an ordered choice, that depends on the neighbours' choices. \\
Spatial Interaction Models & The data are flows between places, such as migration, commuting or trade. \\
Spatial Panel Data & The same areas are observed over several periods. \\
Spatio-Temporal Forecasting & You want forecasts of areas observed over many periods, from their own past and their neighbours'. \\
Spatial Causal Inference & You want the effect of a policy on the areas under it and on the areas around them. \\
Shift-Share Designs & Areas differ in how much a common set of shocks reaches them, through their shares in the shocks. \\
Regionalisation and Spatial Regimes & You want contiguous regions alike inside, or coefficients that differ from one region to another. \\
Spatial Correlation Robust Inference & You want standard errors of a regression that allow for correlation between nearby areas. \\\bottomrule
\end{tabularx}}

\medskip
The panels share the steps that matter in any spatial model: who counts as a neighbour, set by contiguity or by the nearest areas; Moran's I of the residuals, to see whether spatial dependence is left; and maps of the results.

\nmsection{Spatial Lag and Error Models}
Five tabs, from least squares to the spatial Durbin model, as set out by \nmcite{anselin1988}{Anselin (1988)} and \nmcite{lesage2009}{LeSage and Pace (2009)}.

\medskip
{\small
\begin{tabularx}{\linewidth}{@{}>{\raggedright\arraybackslash}p{36mm}>{\raggedright\arraybackslash}p{58mm}X@{}}
\toprule
\textbf{Tab} & \textbf{Spatial term} & \textbf{Use it when} \\\midrule
OLS and LM Tests & None & You first want to know which spatial term the data call for \\
Spatial Lag (SAR) & The neighbours' outcome & Outcomes spill over, as prices or tax rates do \\
Spatial Error (SEM) & The neighbours' errors & Omitted variables shared across borders \\
Lag of X (SLX) & The neighbours' characteristics & Spillovers stay local, with no feedback \\
Spatial Durbin (SDM) & Both the outcome and the characteristics & You are unsure; it nests the other three \\\bottomrule
\end{tabularx}}

\medskip
In SAR and SDM a coefficient is not the effect: Direct and Indirect Effects splits it into the effect on the area itself and on the others. Likelihood Ratio Tests of the SDM test whether it reduces to a simpler model. The Replicate steps set the estimates beside those of \nmcite{bivand2015}{Bivand and Piras (2015)} and of R's spatialreg.

\nmsection{Geographically Weighted Regression}
Coefficients that vary over the map, estimated at each area from its neighbours weighted by distance (\nmcite{brunsdon1996}{Brunsdon, Fotheringham and Charlton 1996}). GWR, Fixed Bandwidth uses one distance, GWR, Adaptive a number of nearest areas, both chosen by AICc (\nmcite{hurvich1998}{Hurvich, Simonoff and Tsai 1998}). Multiscale GWR gives each coefficient its own bandwidth (\nmcite{fotheringham2017}{Fotheringham, Yang and Kang 2017}), so a coefficient that is constant over the map can be told from one that changes over short distances. The actions map the local coefficients and t-values, adjust the local tests for multiple testing (\nmcite{dasilva2016}{da Silva and Fotheringham 2016}), test whether a coefficient varies at all, and set the results beside GWR4 and \nmcite{oshan2019}{Oshan, Li, Kang, Wolf and Fotheringham (2019)} on the Georgia counties.

\nmsection{Spatial Discrete Choice}
Binary and ordered choices that depend on the neighbours' choices. The spatial lag probit is fitted by linearised GMM (\nmcite{klier2008}{Klier and McMillen 2008}), fast on many areas, or by Gibbs sampling (\nmcite{albert1993}{Albert and Chib 1993}; \nmcite{lesage2009}{LeSage and Pace 2009}), which gives the whole posterior. Ordered Probit handles choices with more than two ranked levels, and Logit With a Filter replaces the spatial lag with Moran eigenvectors. As in the spatial lag model, the effect on the probability has a direct and an indirect part. The Replicate step sets the estimates beside \nmcite{lesage2011}{LeSage, Pace, Lam, Campanella and Liu (2011)} on the reopening of New Orleans businesses after Hurricane Katrina.

\nmsection{Spatial Interaction Models}
Flows between places on what each origin sends, what each destination attracts and the distance between them. The four tabs are the unconstrained, production-, attraction- and doubly constrained models of \nmcite{wilson1971}{Wilson (1971)}, fitted as Poisson regressions (\nmcite{flowerdew1982}{Flowerdew and Aitkin 1982}; \nmcite{santossilva2006}{Santos Silva and Tenreyro 2006}) whose origin and destination effects reproduce the known totals; a fifth adds a negative binomial variance. The actions map the flows and the effects, test for overdispersion and for dependence between neighbouring flows, compare the constraints, the forms of distance decay and the standard errors, and predict the flows after distances shrink. The example is the migration between the regions of Austria, and the Replicate step sets the four models beside \nmcite{oshan2016}{Oshan (2016)}'s primer on PySAL's spint.

\nmsection{Spatial Panel Data}
SAR, SEM and SDM with fixed effects of the areas, of the periods or both, removed by the transformation of \nmcite{lee2010}{Lee and Yu (2010)} or by demeaning (\nmcite{elhorst2014}{Elhorst 2014}). The Dynamic Panel adds the outcome's own lag and its neighbours', corrected for the bias the fixed effects cause (\nmcite{yu2008}{Yu, de Jong and Lee 2008}), and its effects come in the short and the long run. The actions run the panel LM tests, compute the standard errors of \nmcite{driscoll1998}{Driscoll and Kraay (1998)}, test for a structural break, and compare the specifications. The Replicate step sets the spatial error estimates beside those of \nmcite{millo2012}{Millo and Piras (2012)} on Munnell's production function for the US states.

\nmsection{Spatio-Temporal Forecasting}
Three models, each fitted on the periods before those held back and scored on them by RMSE and MAE. STAR, the space-time autoregression of \nmcite{pfeifer1980}{Pfeifer and Deutsch (1980)}, forecasts one variable from its own lags and its neighbours'. Spatial VAR, after \nmcite{digiacinto2010}{Di Giacinto (2010)}, forecasts several variables together. Panel AR leaves space out and is the benchmark: a spatial model that does not beat it adds nothing. The actions compare accuracy horizon by horizon with each area's last value, test the residuals with the CD test of \nmcite{pesaran2015}{Pesaran (2015)}, check stationarity, forecast beyond the data and trace space-time impulse responses.

\nmsection{Spatial Causal Inference}
Two designs for a policy whose effects can cross borders.

\medskip
{\small
\begin{tabularx}{\linewidth}{@{}>{\raggedright\arraybackslash}p{26mm}>{\raggedright\arraybackslash}p{60mm}X@{}}
\toprule
\textbf{Tab} & \textbf{What it estimates} & \textbf{Use it when} \\\midrule
Spatial DiD & The direct effect and the spillover from the share of neighbours under the policy (\nmcite{delgado2015}{Delgado and Florax 2015}) & Many areas, some under the policy, observed before and after \\
Synthetic Control & The gap between an area and a weighted mix of comparison areas (\nmcite{abadie2010}{Abadie, Diamond and Hainmueller 2010}), with the penalty of \nmcite{abadie2021}{Abadie and L'Hour (2021)} & One or a few areas under the policy \\\bottomrule
\end{tabularx}}

\medskip
Synthetic Control leaves the neighbours of the areas under the policy out of the comparison areas, so a spillover does not bias the gap, and weighs the predictors equally or by the nested search of Abadie, Diamond and Hainmueller. The examples are the minimum wage panel of US counties of \nmcite{callaway2021}{Callaway and Sant'Anna (2021)} and the cigarette sales of the US states with Proposition 99 in California. The actions check what the design can identify before any outcome and the fit before the policy. A permutation test ranks the gap against random sets of comparison areas. The Replicate step sets the synthetic California beside Tables~1 and~2 of Abadie, Diamond and Hainmueller.

\nmsection{Shift-Share Designs}
Shift-Share IV estimates the effect of a change, such as the growth of imports in a local labour market, instrumented by each area's shares in a set of shocks times the shocks; Reduced Form regresses the outcome on that shift-share regressor itself. Both estimate across the shocks, which \nmcite{borusyak2022}{Borusyak, Hull and Jaravel (2022)} show gives the same coefficient, with a standard error valid when the shocks are as good as randomly assigned, however the areas' residuals are correlated through similar shares (\nmcite{adao2019}{Ad\~{a}o, Koles\'{a}r and Morales 2019}). Each tab prints that standard error, clustered by groups of shocks, the confidence interval with the null imposed and the effective number of shocks. The example is the China shock of \nmcite{autor2013}{Autor, Dorn and Hanson (2013)} across the US commuting zones. The actions check the concentration and balance of the shocks, set the standard error beside one clustered by region, show which shocks drive the estimate (\nmcite{goldsmith2020}{Goldsmith-Pinkham, Sorkin and Swift 2020}), and recompute it shock by shock. The Replicate step sets the seven columns of Table~4 of Borusyak, Hull and Jaravel beside Julia's.

\nmsection{Regionalisation and Spatial Regimes}
SKATER Regions groups the areas into a chosen number of contiguous regions, alike inside in the variables chosen, by cutting a minimum spanning tree of neighbouring areas (\nmcite{assuncao2006}{Assun\c{c}\~{a}o, Neves, C\^{a}mara and Freitas 2006}). Regimes, OLS gives each regime its own intercept and slopes, with the test of \nmcite{chow1960}{Chow (1960)} that the split is needed. Regimes, Spatial Lag adds a spatial lag of the outcome with one rho shared by the regimes, since ignoring spatial dependence invalidates the usual tests of equal coefficients (\nmcite{anselin1990}{Anselin 1990}). The regimes come from a column, such as a city and its suburbs, or from SKATER; regions found from the data make the Chow test too optimistic, and the Bootstrap Chow Test repeats the search on data with no regimes. The actions test which coefficients differ, choose the number of regions, and compare the regimes with a single model and with geographically weighted regression. The Replicate step sets the regime spatial lag model beside PySAL's spreg on the Baltimore house prices of \nmcite{dubin1992}{Dubin (1992)}.

\nmsection{Spatial Correlation Robust Inference}
Least squares, with the coefficients' standard errors corrected for correlation between the areas' errors rather than the dependence modelled. Conley Spatial HAC allows for correlation between areas closer than a cutoff, weighted by a kernel of the distance (\nmcite{conley1999}{Conley 1999}), as \nmcite{newey1987}{Newey and West (1987)} do over time. Cluster-Robust allows for any correlation inside groups of areas, given many groups (\nmcite{cameron2015}{Cameron and Miller 2015}). Driscoll--Kraay, for panels with many periods, allows for any correlation between areas in the same period (\nmcite{driscoll1998}{Driscoll and Kraay 1998}). The actions show the residuals' correlation by distance, the standard errors across cutoffs, kernels, cluster levels and lags, and all the estimators side by side. The Replicate steps set the standard errors beside the test data of \nmcite{petersen2009}{Petersen (2009)}, R's plm and fixest.

\nmsection{References}
{\footnotesize\sloppy
\nmref{abadie2010}{Abadie, A., Diamond, A. and Hainmueller, J. (2010). Synthetic control methods for comparative case studies: estimating the effect of California's tobacco control program. \emph{Journal of the American Statistical Association} 105(490), 493--505. \nmdoi{10.1198/jasa.2009.ap08746}}
\nmref{abadie2021}{Abadie, A. and L'Hour, J. (2021). A penalized synthetic control estimator for disaggregated data. \emph{Journal of the American Statistical Association} 116(536), 1817--1834. \nmdoi{10.1080/01621459.2021.1971535}}
\nmref{adao2019}{Ad\~{a}o, R., Koles\'{a}r, M. and Morales, E. (2019). Shift-share designs: theory and inference. \emph{The Quarterly Journal of Economics} 134(4), 1949--2010. \nmdoi{10.1093/qje/qjz025}}
\nmref{albert1993}{Albert, J. H. and Chib, S. (1993). Bayesian analysis of binary and polychotomous response data. \emph{Journal of the American Statistical Association} 88(422), 669--679. \nmdoi{10.1080/01621459.1993.10476321}}
\nmref{anselin1988}{Anselin, L. (1988). \emph{Spatial Econometrics: Methods and Models}. Springer Netherlands. \nmdoi{10.1007/978-94-015-7799-1}}
\nmref{anselin1990}{Anselin, L. (1990). Spatial dependence and spatial structural instability in applied regression analysis. \emph{Journal of Regional Science} 30(2), 185--207. \nmdoi{10.1111/j.1467-9787.1990.tb00092.x}}
\nmref{assuncao2006}{Assun\c{c}\~{a}o, R. M., Neves, M. C., C\^{a}mara, G. and Freitas, C. da C. (2006). Efficient regionalization techniques for socio-economic geographical units using minimum spanning trees. \emph{International Journal of Geographical Information Science} 20(7), 797--811. \nmdoi{10.1080/13658810600665111}}
\nmref{autor2013}{Autor, D. H., Dorn, D. and Hanson, G. H. (2013). The China syndrome: local labor market effects of import competition in the United States. \emph{American Economic Review} 103(6), 2121--2168. \nmdoi{10.1257/aer.103.6.2121}}
\nmref{bivand2015}{Bivand, R. and Piras, G. (2015). Comparing implementations of estimation methods for spatial econometrics. \emph{Journal of Statistical Software} 63(18). \nmdoi{10.18637/jss.v063.i18}}
\nmref{borusyak2022}{Borusyak, K., Hull, P. and Jaravel, X. (2022). Quasi-experimental shift-share research designs. \emph{The Review of Economic Studies} 89(1), 181--213. \nmdoi{10.1093/restud/rdab030}}
\nmref{brunsdon1996}{Brunsdon, C., Fotheringham, A. S. and Charlton, M. E. (1996). Geographically weighted regression: a method for exploring spatial nonstationarity. \emph{Geographical Analysis} 28(4), 281--298. \nmdoi{10.1111/j.1538-4632.1996.tb00936.x}}
\nmref{callaway2021}{Callaway, B. and Sant'Anna, P. H. C. (2021). Difference-in-differences with multiple time periods. \emph{Journal of Econometrics} 225(2), 200--230. \nmdoi{10.1016/j.jeconom.2020.12.001}}
\nmref{cameron2015}{Cameron, A. C. and Miller, D. L. (2015). A practitioner's guide to cluster-robust inference. \emph{Journal of Human Resources} 50(2), 317--372. \nmdoi{10.3368/jhr.50.2.317}}
\nmref{chow1960}{Chow, G. C. (1960). Tests of equality between sets of coefficients in two linear regressions. \emph{Econometrica} 28(3), 591--605. \nmdoi{10.2307/1910133}}
\nmref{conley1999}{Conley, T. G. (1999). GMM estimation with cross sectional dependence. \emph{Journal of Econometrics} 92(1), 1--45. \nmdoi{10.1016/S0304-4076(98)00084-0}}
\nmref{dasilva2016}{da Silva, A. R. and Fotheringham, A. S. (2016). The multiple testing issue in geographically weighted regression. \emph{Geographical Analysis} 48(3), 233--247. \nmdoi{10.1111/gean.12084}}
\nmref{delgado2015}{Delgado, M. S. and Florax, R. J. G. M. (2015). Difference-in-differences techniques for spatial data: local autocorrelation and spatial interaction. \emph{Economics Letters} 137, 123--126. \nmdoi{10.1016/j.econlet.2015.10.035}}
\nmref{digiacinto2010}{Di Giacinto, V. (2010). On vector autoregressive modeling in space and time. \emph{Journal of Geographical Systems} 12(2), 125--154. \nmdoi{10.1007/s10109-010-0116-6}}
\nmref{driscoll1998}{Driscoll, J. C. and Kraay, A. C. (1998). Consistent covariance matrix estimation with spatially dependent panel data. \emph{The Review of Economics and Statistics} 80(4), 549--560. \nmdoi{10.1162/003465398557825}}
\nmref{dubin1992}{Dubin, R. A. (1992). Spatial autocorrelation and neighborhood quality. \emph{Regional Science and Urban Economics} 22(3), 433--452. \nmdoi{10.1016/0166-0462(92)90038-3}}
\nmref{elhorst2014}{Elhorst, J. P. (2014). \emph{Spatial Econometrics: From Cross-Sectional Data to Spatial Panels}. Springer Berlin Heidelberg. \nmdoi{10.1007/978-3-642-40340-8}}
\nmref{flowerdew1982}{Flowerdew, R. and Aitkin, M. (1982). A method of fitting the gravity model based on the Poisson distribution. \emph{Journal of Regional Science} 22(2), 191--202. \nmdoi{10.1111/j.1467-9787.1982.tb00744.x}}
\nmref{fotheringham2017}{Fotheringham, A. S., Yang, W. and Kang, W. (2017). Multiscale geographically weighted regression (MGWR). \emph{Annals of the American Association of Geographers} 107(6), 1247--1265. \nmdoi{10.1080/24694452.2017.1352480}}
\nmref{goldsmith2020}{Goldsmith-Pinkham, P., Sorkin, I. and Swift, H. (2020). Bartik instruments: what, when, why, and how. \emph{American Economic Review} 110(8), 2586--2624. \nmdoi{10.1257/aer.20181047}}
\nmref{hurvich1998}{Hurvich, C. M., Simonoff, J. S. and Tsai, C.-L. (1998). Smoothing parameter selection in nonparametric regression using an improved Akaike information criterion. \emph{Journal of the Royal Statistical Society Series B} 60(2), 271--293. \nmdoi{10.1111/1467-9868.00125}}
\nmref{klier2008}{Klier, T. and McMillen, D. P. (2008). Clustering of auto supplier plants in the United States: generalized method of moments spatial logit for large samples. \emph{Journal of Business \& Economic Statistics} 26(4), 460--471. \nmdoi{10.1198/073500107000000188}}
\nmref{lee2010}{Lee, L.-f. and Yu, J. (2010). Estimation of spatial autoregressive panel data models with fixed effects. \emph{Journal of Econometrics} 154(2), 165--185. \nmdoi{10.1016/j.jeconom.2009.08.001}}
\nmref{lesage2009}{LeSage, J. and Pace, R. K. (2009). \emph{Introduction to Spatial Econometrics}. Chapman and Hall/CRC. \nmdoi{10.1201/9781420064254}}
\nmref{lesage2011}{LeSage, J. P., Pace, R. K., Lam, N., Campanella, R. and Liu, X. (2011). New Orleans business recovery in the aftermath of Hurricane Katrina. \emph{Journal of the Royal Statistical Society Series A} 174(4), 1007--1027. \nmdoi{10.1111/j.1467-985X.2011.00712.x}}
\nmref{millo2012}{Millo, G. and Piras, G. (2012). splm: spatial panel data models in R. \emph{Journal of Statistical Software} 47(1), 1--38. \nmdoi{10.18637/jss.v047.i01}}
\nmref{newey1987}{Newey, W. K. and West, K. D. (1987). A simple, positive semi-definite, heteroskedasticity and autocorrelation consistent covariance matrix. \emph{Econometrica} 55(3), 703--708. \nmdoi{10.2307/1913610}}
\nmref{oshan2016}{Oshan, T. M. (2016). A primer for working with the Spatial Interaction modeling (SpInt) module in the Python Spatial Analysis Library (PySAL). \emph{REGION} 3(2), R11--R23. \nmdoi{10.18335/region.v3i2.175}}
\nmref{oshan2019}{Oshan, T. M., Li, Z., Kang, W., Wolf, L. J. and Fotheringham, A. S. (2019). mgwr: a Python implementation of multiscale geographically weighted regression for investigating process spatial heterogeneity and scale. \emph{ISPRS International Journal of Geo-Information} 8(6), 269. \nmdoi{10.3390/ijgi8060269}}
\nmref{pesaran2015}{Pesaran, M. H. (2015). Testing weak cross-sectional dependence in large panels. \emph{Econometric Reviews} 34(6--10), 1089--1117. \nmdoi{10.1080/07474938.2014.956623}}
\nmref{petersen2009}{Petersen, M. A. (2009). Estimating standard errors in finance panel data sets: comparing approaches. \emph{The Review of Financial Studies} 22(1), 435--480. \nmdoi{10.1093/rfs/hhn053}}
\nmref{pfeifer1980}{Pfeifer, P. E. and Deutsch, S. J. (1980). A three-stage iterative procedure for space-time modeling. \emph{Technometrics} 22(1), 35--47. \nmdoi{10.2307/1268381}}
\nmref{santossilva2006}{Santos Silva, J. M. C. and Tenreyro, S. (2006). The log of gravity. \emph{The Review of Economics and Statistics} 88(4), 641--658. \nmdoi{10.1162/rest.88.4.641}}
\nmref{wilson1971}{Wilson, A. G. (1971). A family of spatial interaction models, and associated developments. \emph{Environment and Planning A} 3(1), 1--32. \nmdoi{10.1068/a030001}}
\nmref{yu2008}{Yu, J., de Jong, R. and Lee, L.-f. (2008). Quasi-maximum likelihood estimators for spatial dynamic panel data with fixed effects when both n and T are large. \emph{Journal of Econometrics} 146(1), 118--134. \nmdoi{10.1016/j.jeconom.2008.08.002}}
}

\nmcredit
\end{document}
